\documentclass[11pt,a4paper]{article}
\usepackage[margin=1in]{geometry}
\usepackage{amsmath,amssymb,amsthm}
\usepackage{booktabs,tabularx,array}
\usepackage{algorithm}
\usepackage{algpseudocode}
\usepackage{enumitem}
\usepackage{microtype}
\usepackage{xcolor}
\usepackage[hidelinks]{hyperref}
\usepackage{fancyhdr}
\usepackage{caption}
\pagestyle{fancy}
\fancyhf{}
\fancyhead[L]{\small Provenance-Aware Symbolic Rule Fusion}
\fancyhead[R]{\small Reinforcement Learning Auditability}
\fancyfoot[C]{\thepage}
\newtheorem{definition}{Definition}
\newtheorem{remark}{Remark}
\newcommand{\E}{\mathbb{E}}
\newcommand{\R}{\mathcal{R}}
\newcommand{\X}{\mathcal{X}}
\newcommand{\K}{\mathcal{K}}

\title{From Replayable Traces to Provenance-Aware Rule Fusion in Reinforcement Learning\\
Evidence, Failure Modes, and Measurement Boundaries}
\author{Anonymous Authors}
\date{September 2026}

\begin{document}
\maketitle

\begin{abstract}
Replay integrity, symbolic rule overlap, behavioral agreement, and policy composition are distinct properties of reinforcement-learning artifacts. We report a bounded end-to-end study that measures them separately. Eight REINFORCE agents per task were trained on CartPole-v1 and Acrobot-v1 with one frozen four-bin state symbolizer shared within each task. On 100 matched reset seeds, provenance-aware rule fusion obtained mean returns of 205.96 on CartPole and $-443.37$ on Acrobot, compared with 163.65 and $-500.00$ for the best single actors; the paired bootstrap intervals were $[23.98,61.93]$ and $[39.32,75.15]$. These gains do not establish optimal arbitration. Exact rule Jaccard for two CartPole actors was 0.5833, whereas agreement on 10,000 fresh raw states was only 51.36\% (95\% interval 50.48--52.26\%). Acrobot arbitration was unexpectedly sensitive: random selection among available source rules averaged $-198.82$ across five selector streams, while the tested confidence-ranked arbiter reached $-443.37$ and uniform random actions averaged $-498.71$. Fitted-Q generalised policy improvement failed diagnostic quality checks and is therefore not treated as a positive comparator. The evidence supports bounded replay and provenance capabilities under the evaluated protocol, not universal auditability or a causal explanation of arbitrary neural decisions.
\end{abstract}

\section{Introduction}
\label{sec:intro}

Reinforcement-learning systems are usually evaluated by return, but a useful scientific account also asks what was recorded, what can be replayed, and which parts of a decision can be inspected. These questions are easy to conflate. A trajectory can be replayable without exposing the reason for an action. Two agents can share a symbolic production without choosing the same action on raw states. A policy merger can improve return while its arbitration rule remains poorly understood.

We therefore study auditability as a collection of separable properties rather than a single score. The work proceeds from a simple starting point---actors produce trajectories and returns---to increasingly stronger objects: state symbols, admitted rules, provenance records, replay predicates, behavioral tests, and policy composition. The sequence is intentional. Several early approaches looked plausible but failed because they tested the wrong object, used an invalid environment boundary, or confused a diagnostic signal with an identification-bearing result.

The central claim is bounded: under a documented two-task protocol, a workflow can expose replayable execution records, compose empirical state-to-action rules with visible source provenance, and report conflicts and blind spots. We do not claim that the resulting grammar is a causal neural explanation, that exact overlap is behavioral consensus, or that arbitrary RL training becomes fully auditable.

\paragraph{Contributions.}
\begin{enumerate}[leftmargin=*,itemsep=2pt]
  \item We formalize shared quantile symbolization, rule admission, provenance-aware arbitration, blind-spot fallback, and replay checks.
  \item We report multi-seed fusion results together with a fresh-state behavioral-agreement test that directly challenges a tempting overlap interpretation.
  \item We retain negative results: lossy grammar prototypes, invalid replay audits, strict-threshold starvation, compression non-gains, arbitration surprises, and failed fitted-Q diagnostics.
  \item We state the strongest supported auditability boundary and separate operational integrity from semantic and causal claims.
\end{enumerate}

\section{Related Work and Positioning}
Programmatically interpretable reinforcement learning seeks policies represented by programs subject to performance constraints \cite{verma2018pirl}. Successor features and generalised policy improvement compose value knowledge across policies or tasks \cite{barreto2017sf,thakoor2022gpi}. Sequence grammars motivate lossless descriptions of observed strings \cite{nevill1997sequitur}, while work on emergent communication cautions that correlation and causal use must be separated \cite{lowe2019pitfalls}. Reproducibility work in deep RL motivates explicit seeding and deterministic checks \cite{nagarajan2018det}.

Our setting is narrower than a general interpretable-RL claim. We fuse empirical state-to-action rules extracted from multiple seeds of the same control task, attach source and decision provenance, and test environment replay. Value-based composition is included as a diagnostic comparator, not assumed valid in advance. The contribution is consequently an evidence-bounded instrument and composition study rather than a claim of a new universal policy representation.

\section{Methods}
\label{sec:methods}

\subsection{Starting point: actors and calibration}
We trained 64-unit REINFORCE actors on CartPole-v1 and Acrobot-v1 for 300 episodes per seed and arm. The seed set was
\[
 S=\{11,29,43,71,101,149,211,307\}.
\]
For each task, 24 random-policy calibration episodes per seed were pooled. For observation dimension $d$ and four bins, the frozen empirical edges were
\begin{equation}
\eta_{d,k}=\operatorname{Quantile}_{k/4}(\{x_d:x\in\mathcal C\}),\qquad k\in\{1,2,3\},
\label{eq:edges}
\end{equation}
and the symbolizer was
\begin{equation}
\phi(x)_d=\sum_{k=1}^{3}\mathbf{1}[x_d\geq\eta_{d,k}].
\label{eq:symbolizer}
\end{equation}
The common calibration removes a seed-specific discretization artifact from overlap statistics; it does not make different rules semantically equivalent.

\subsection{What failed before the validated pipeline}
The first implementation stage exposed why an audit claim must follow the actual runtime boundary. An early example passed an unsupported constructor argument, reset the environment twice while attempting to seed it, and treated termination as the only episode boundary even though the environment also reports truncation. Its sequence inducer greedily replaced repeated bigrams, truncated a start production, and inferred provenance through substring matching. The accompanying auditor searched production prefixes but did not re-execute the policy or the environment. These outputs were therefore diagnostic artifacts, not evidence of replayability.

A second failure came from schema design. A grammar sidecar field was included in an update-ledger equality check even though it did not affect optimization. The resulting mismatch falsely suggested that the baseline and audit-only arms diverged. Removing the sidecar from the equality core repaired the comparison while preserving the incident as a warning: an audit schema can manufacture a false scientific difference.

Two further probes constrained the design. Applying high-confidence lookup constraints online reduced one CartPole actor's held-out return from 500 to 121.3 at confidence 0.90 while covering roughly 0.4\% of steps. A strict confidence-0.90 merge probe produced 100\% blind spots. These failures motivated an evaluated rule-fusion policy with explicit fallback and a primary threshold of support 8 and confidence 0.70, rather than treating high confidence as automatically safe.

Finally, a lossless pair-substitution grammar reconstructed observed traces but did not beat ordinary gzip in the storage comparison. A two-part description-length ratio for one seed rose from about 0.513 to 0.575 across rounds. We therefore treat trace grammar and state-action rule banks as different objects and make no compression claim.

\subsection{Rule admission and fusion}
For a symbolized condition $c$ and action $a$, let $N(c,a)$ count training visits. The rule statistics are
\begin{equation}
\operatorname{supp}(c,a)=N(c,a),\qquad
\operatorname{conf}(c,a)=\frac{N(c,a)}{\sum_{a'}N(c,a')}.
\label{eq:statistics}
\end{equation}
A production enters the bank $\R_s$ only when
\begin{equation}
\rho=(c\mapsto a)\in\R_s
\quad\Longleftrightarrow\quad
\operatorname{supp}(c,a)\geq 8\;\wedge\;
\operatorname{conf}(c,a)\geq 0.70.
\label{eq:admission}
\end{equation}
At observation $x$, the available candidates are
\begin{equation}
\K(x)=\{\rho\in\bigcup_{s\in S}\R_s:\operatorname{cond}(\rho)=\phi(x)\}.
\label{eq:candidates}
\end{equation}
The default arbiter selects lexicographically by confidence, support, mean reward, and then source priority:
\begin{equation}
\rho^*(x)=\arg\max_{\rho\in\K(x)}
  (\operatorname{conf}(\rho),\operatorname{supp}(\rho),\bar r(\rho),-s(\rho)).
\label{eq:arbiter}
\end{equation}
The fused policy is
\begin{equation}
\pi^{\mathrm{fus}}(x)=
\begin{cases}
\operatorname{act}(\rho^*(x)),&\K(x)\neq\varnothing,\\
\pi^{\mathrm{fb}}(x),&\K(x)=\varnothing,
\end{cases}
\label{eq:fusion}
\end{equation}
where the fallback actor is selected from training-return metadata only.

\begin{algorithm}[t]
\caption{Provenance-aware rule-fusion decision}
\label{alg:fusion}
\begin{algorithmic}[1]
\Require observation $x$, symbolizer $\phi$, rule banks $\{\R_s\}$, fallback $\pi^{\mathrm{fb}}$
\Ensure action $a$ and provenance record $\nu$
\State $c\gets\phi(x)$; $\K\gets\{\rho\in\cup_s\R_s:\operatorname{cond}(\rho)=c\}$
\If{$\K=\varnothing$}
  \State $a\gets\pi^{\mathrm{fb}}(x)$; mark $\nu$ as a blind spot
\ElsIf{all candidates have one action}
  \State choose that action and mark no conflict
\Else
  \State choose $\rho^*$ by Eq.~\eqref{eq:arbiter}; attach its source and mark conflict
\EndIf
\State append $(c,a,\nu)$ to the decision log and update the digest
\State \Return $(a,\nu)$
\end{algorithmic}
\end{algorithm}

\subsection{Integrity, replay, and behavioral estimands}
For payload $p_i$ and previous digest $h_{i-1}$, the chain predicate is
\begin{equation}
\operatorname{ChainValid}\iff\forall i,\quad h_i=H(h_{i-1}\Vert p_i),
\label{eq:chain}
\end{equation}
and a replayed transition must satisfy
\begin{equation}
\operatorname{ReplayOK}_t\iff(x'_{t+1},r'_t,\operatorname{term}'_t,\operatorname{trunc}'_t)
 =(x_{t+1},r_t,\operatorname{term}_t,\operatorname{trunc}_t).
\label{eq:replay}
\end{equation}
These predicates identify consistency of recorded artifacts under the tested stack; they do not prove logger honesty or cross-hardware equivalence.

Exact rule Jaccard is a syntactic statistic. Experiment B instead measured fresh-state argmax agreement:
\begin{equation}
\alpha=\Pr_{x\sim\mathcal D}
 [\arg\max_a\pi_{11}(a\mid x)=\arg\max_a\pi_{29}(a\mid x)].
\label{eq:agreement}
\end{equation}
For $M=100$ matched reset seeds, mean return and paired difference were
\begin{equation}
\hat\mu(\pi)=\frac1M\sum_{m=1}^{M}R(\pi;z_m),\qquad
\Delta=\hat\mu(\pi)-\hat\mu(\pi_{\mathrm{best}}),
\label{eq:return}
\end{equation}
with percentile bootstrap intervals on paired episode-return differences.

\section{Results}
\label{sec:results}

\subsection{Primary returns}
\begin{table}[t]
\centering
\caption{Mean return under 100 matched resets. Intervals are paired bootstrap intervals versus the best single actor.}
\label{tab:returns}
\small
\begin{tabular}{@{}lrrrr@{}}
\toprule
Policy & CartPole mean & CartPole $\Delta$ & Acrobot mean & Acrobot $\Delta$ \\
\midrule
Best single actor & 163.65 & --- & $-500.00$ & --- \\
Actor-logit ensemble & 500.00 & $+336.35$ & $-500.00$ & $0.00$ \\
Rule fusion & 205.96 & $+42.31$ & $-443.37$ & $+56.63$ \\
\bottomrule
\end{tabular}
\end{table}

Rule fusion improved mean return relative to the best single actor on both tasks (Table~\ref{tab:returns}). The CartPole scale saturates at 500 for some policies. On Acrobot, fusion covered 95.70\% of decision steps, encountered conflicts on 89.24\%, and had 4.30\% blind spots; the median remained $-500$, while 39 of 100 episodes ended above the floor. These distributional facts prevent the mean difference from being read as uniform improvement.

\subsection{Overlap is not behavioral consensus}
Across eight seeds, mean pairwise exact grammar Jaccard was 0.5572 for CartPole and 0.0956 for Acrobot. Seeds 11 and 29 had CartPole Jaccard 0.5833 with no action conflict among 77 shared conditions. Yet on 10,000 fresh raw states their pooled action agreement was 51.38\%, and the episode-cluster balanced estimate was 51.36\% (95\% interval 50.48--52.26\%). Acrobot had Jaccard 0.0496 and 77 conflicts among 98 shared conditions. Thus Jaccard alone cannot serve as evidence of a shared behavioral core.

\subsection{Arbitration sensitivity and failed value composition}
The confidence/support/reward arbiter scored $-443.37$ on Acrobot; confidence-only scored $-444.49$, while mean-reward-only, support-first, and conflict deferral to the actor ensemble scored $-500$. Random selection among available source rules scored $-187.02$ on one stream and $-198.82$ averaged over five streams. Uniform random actions averaged $-498.71$. The state-conditioned random-rule effect is therefore not explained by unconditioned action noise, but its mechanism remains unknown. It is a finding requiring follow-up, not a general arbitration recommendation.

FQE-GPI achieved mean return 9.34 on CartPole and $-499.57$ on Acrobot. Its CartPole actions agreed with the actor-logit ensemble on only 35.27\% of 31,563 held-out states; on Acrobot it selected action 0 throughout a 50,000-state diagnostic sample. On 200 Acrobot states, the highest continuation-return action matched fusion only 22.5\% of the time (95\% interval 16.5--28.5\%). These diagnostics fail quality checks, so they are not positive evidence against or for the primary fusion result.

\subsection{Trace-scale audit evidence}
Hash-chain, update-ledger, provenance-resolution, and environment-replay summaries passed under the audited setup. The checked records covered 1,725,655 CartPole rows and 2,544,513 Acrobot rows, with 1,122,403 and 1,742,570 replayed transitions. These are operational and protocol-level checks. They do not identify every training cause, prove adversarial integrity, or establish semantic adequacy.

\section{Discussion}
\label{sec:discussion}
The positive result is an existence result for a bounded instrument: source rules can be admitted, selected, logged with provenance, and replay-checked while producing a measurable policy. The negative results are equally important. A syntactic overlap measure overstates behavioral agreement in one direct test. Thresholds can starve coverage. Online rule constraints can damage return. The declared arbiter is not validated by the strongest Acrobot ablation. A value-composition baseline collapses diagnostic checks. Together these observations constrain the scientific claim rather than merely complicating implementation.

The strongest alternative explanation for the Acrobot gain is that rule availability and action-frequency structure, rather than confidence, determine the outcome. The present artifacts do not identify that mechanism. Another explanation for the CartPole gain is scale saturation and the weakness of the best-single fallback, not generally superior composition. The paired design supports a within-reset comparison, but not broad transfer across tasks or hardware.

The auditability layers should remain separate: replay validates recorded transitions; provenance validates references and decision history; behavioral agreement measures action coincidence on fresh states; and causal identification would require interventions that are not part of the control-task fusion experiment. A passing layer cannot substitute for another.

\section{Limitations and Reproducibility}
The evidence covers two classic-control tasks, eight seeds per task, one shared quantizer per task, and a particular software stack. CartPole has a ceiling and Acrobot has a floor, limiting return resolution. Fresh behavioral agreement was directly measured for one CartPole pair rather than every pair. The random-rule mechanism is unresolved. FQE-GPI is not a credible positive comparator in this study. Cross-hardware replay and adversarial logger checks were not performed. The reviewed compact artifacts preserve summaries and hashes, while the full traces and checkpoints are external to the compact package.

\section{Conclusion}
Under the documented protocol, provenance-aware state-action rule fusion supports inspectable conflicts, blind spots, source references, and replay checks, and it improves mean return relative to the selected best single actor on the two evaluated tasks. The same evidence shows that grammar overlap is not behavioral consensus and that arbitration and value-composition claims remain unsettled. The defensible conclusion is bounded auditability of specified artifacts, not universal auditability or complete causal explanation of reinforcement learning.

\begin{thebibliography}{9}
\bibitem{verma2018pirl} A. Verma, V. Murali, R. Singh, P. Kohli, and S. Chaudhuri. Programmatically Interpretable Reinforcement Learning. In \emph{Proceedings of ICML}, 2018.
\bibitem{barreto2017sf} A. Barreto et al. Successor Features for Transfer in Reinforcement Learning. In \emph{Advances in Neural Information Processing Systems}, 2017.
\bibitem{thakoor2022gpi} S. Thakoor et al. Generalised Policy Improvement with Geometric Policy Composition. In \emph{Proceedings of ICML}, 2022.
\bibitem{nevill1997sequitur} C. G. Nevill-Manning and I. H. Witten. Identifying Hierarchical Structure in Sequences: A Linear-Time Algorithm. \emph{Journal of Artificial Intelligence Research}, 7:67--82, 1997.
\bibitem{lowe2019pitfalls} R. Lowe et al. On the Pitfalls of Measuring Emergent Communication. In \emph{Proceedings of AAMAS}, 2019.
\bibitem{nagarajan2018det} P. Nagarajan, G. Warnell, and P. Stone. Deterministic Implementations for Reproducibility in Deep Reinforcement Learning. arXiv:1809.05676, 2018.
\end{thebibliography}
\end{document}


